% =====================================================================
%  Chapter 33 — What a Lunar Station Can and Cannot Be Made Responsible For
% =====================================================================

\chapter{What a Lunar Station Can and Cannot Be Made Responsible For}
\label{ch:close}

The purpose of this final chapter is to make the proposal accountable. A
capability that cannot be given a defined job cannot be held to a standard,
and a claim that cannot be falsified cannot be tested. So this chapter
states the job, states what the station is explicitly not for, and states
the conditions under which the argument is wrong.

\section{What it can be made responsible for}

Four jobs, each of which the preceding chapters have quantified.

\textbf{A resolved shape, on request, for a named target.} A shape model
without a thermal model and without an albedo assumption, for a target that
the ground cannot resolve and that no space observatory is assigned to
resolve. The matched baseline is known from the target's distance; the
aperture, from the target's angular size; and the turnaround, from the
Moon's rotation, which delivers the required orthogonal baseline orientation
in 6.83 days.

\textbf{A range, for a close approach.} A range to a non-cooperative target,
produced optically, with no radar and no cooperative transponder. Outside the
Earth-Moon baseline, by parallax, to a few metres at the distances where
close encounters happen; inside it, directly from the Moon's ephemeris, to
the centimetre.

\textbf{A presence at a close encounter}, and therefore the measurement of
the spin-state change that reveals the mass distribution and with it the
density. This is the only route to the second factor of three in the energy
relation for the many objects that will never receive a spacecraft.

\textbf{An alert path, from the risk list to an instrument.} Not a faster
link; a link at all. The documented failure of the 2024~YR4 episode was
that the observers of the object were not told it had reached Torino
Level 1, and the 2025~FA22 exercise showed the mechanism to be
operationally simple. This job should be done whether or not the station is
built.

\section{What it cannot be made responsible for}

\begin{table}[!htbp]
\centering
\small
\caption{The exclusions, with the reason for each.}
\label{tab:exclusions}
\begin{tabularx}{\textwidth}{@{}P{1.41in}P{1.15in}Y@{}}
\toprule
\rowcolor{panel}
\textbf{Not for} & \textbf{Because} & \textbf{What covers it instead} \\
\midrule
Discovery & it is not a survey, and the 140-metre class is essentially
complete in discovery & the existing survey network \\
\rowcolor{panel}
The thermal channel & the sunlit ground is brighter than the target at
every distance of interest, and L1 is already funded & NEO Surveyor and
NEOMIR \\
A census of the catalogue & 37 campaigns a year against 342 arrivals and
11,725 known & nothing; this is a triage service \\
\rowcolor{panel}
Radar's job & the fourth-power wall and the radar-crossing selection
restrict it to dozens of objects & bistatic radar, and it is indispensable \\
Deflection & outside the scope of this book entirely & other programmes \\
\bottomrule
\end{tabularx}
\end{table}

\tabref{tab:exclusions} is the honest boundary. The two exclusions that
matter most are the census and the thermal channel, because both would be
tempting to claim.

\section{What would falsify this argument}

Four measurements would break it, and they are stated in advance so that
they can be used.

\begin{enumerate}
\item \textbf{A published lunar shape-imaging baseline requirement that is
kilometre-scale.} If a peer-reviewed treatment of a lunar optical
interferometer for asteroid imaging exists, and it specifies a baseline of
hundreds of metres or less with four samples across the diameter, then
Chapter 20's requirement is wrong and the entire lunar case collapses. This
is the most likely single point of failure, and it is checkable by any
reader with access to the literature.
\item \textbf{A published Earth-Moon two-station ranging treatment for a
non-cooperative near-Earth target.} If such a treatment exists, the book's
second new result is prior art, and the contribution reduces to the error
budgets and the rule.
\item \textbf{A measured fringe visibility for a resolved small body at a
matched baseline that is too low to image.} The book uses a uniform-disc
model for the visibility and says so. A real resolved body with limb
darkening and surface contrast could have a visibility so low that the
imaging is impossible at any signal-to-noise. This would not invalidate the
requirement, which is set by the sampling condition, but it would change
the cost.
\item \textbf{A close approach at which the density term turns out to be
smaller than the astrometric term.} The 2029 Apophis encounter is the test
case: the published sensitivity table gives the density a factor of 93 and
the astrometric standard deviation is 399 kilometres against the density's
379 kilometres of displacement. If the astrometric term is in fact smaller
than the physical terms in the real orbit solution, the rule will decline
Apophis too, and the proposal loses its best case.
\end{enumerate}

\section{The closing statement}

There are more than forty thousand near-Earth objects in the catalogue, and
about eleven thousand of them are large enough to matter. Of those, about
one in seven has a spectrum, fewer than one in a hundred has been imaged
with radar, and about nine thousand nine hundred and sixty-six have no
physical account at all.

That is the condition this book was written about, and it is a condition of
\emph{accounting}, not of detection. The dangerous objects have been found.
They have not been sized.

The argument of the book is that the missing measurement is a resolved
shape, and that the measurement can be closed with a resolved shape, a
close-approach mass and a link from the risk list to an instrument. It is
that the lunar surface is the only site where the kilometre-scale stable
baseline that shape imaging requires can be built, and the only site where
a presence at a close encounter and a long observing window come together in
one platform.

And it is that the platform is worth exactly as much as the physical term
is worth, which is nothing at seven years and a great deal at thirty.

\begin{keytakeaways}
\item The station can be made responsible for four jobs: a resolved shape on
request, a range at close approach, a presence at the encounter that reveals
the density, and an alert path that does not exist today.
\item It cannot be made responsible for discovery, for the thermal channel,
for a census of the catalogue, or for deflection.
\item Four measurements would falsify the argument, and the most likely one
is a published lunar shape-imaging design at a shorter baseline than the
requirement.
\item The dangerous objects have been found. They have not been sized, and
that is the condition the book was written about.
\end{keytakeaways}
